\chapter{How to Improve Your Posture}

\section*{Definition and Overview}
Posture is the position in which the body is held when sitting, standing, and moving. Good posture keeps the body's joints, muscles, and spine in natural alignment, reducing strain and preventing pain. Poor posture, common with long hours at desks, phones, and computers, can cause neck, shoulder, and back pain, headaches, and fatigue. Improving posture is not about rigidity or "standing up straight" all day, but about building better habits, strengthening supporting muscles, and taking regular movement breaks. Small, consistent changes produce lasting improvement.

\section*{Epidemiology}
Posture-related pain is extremely common. Around four in five adults experience low back pain at some point, and neck and shoulder pain associated with desk work and device use has increased substantially. "Tech neck" from looking down at phones is now seen in adolescents and young adults. Prolonged sitting, poor workstations, and weak core muscles are major contributors. Posture improvement programs, ergonomic changes, and exercise are effective for the majority of people with postural pain.

\section*{Aetiology and Pathophysiology}
\begin{itemize}
    \item \textbf{Prolonged sitting:} sustained positions load the spine and shorten or weaken muscles
    \item \textbf{Device use:} looking down at phones and laptops rounds the shoulders and flexes the neck
    \item \textbf{Muscle imbalance:} weak core and gluteal muscles with tight chest and hip muscles pull the body out of alignment
    \item \textbf{Poor workstation setup:} screens, chairs, and desks at the wrong height encourage slumping
    \item \textbf{Repetitive activities:} lifting, driving, and caring roles load the body asymmetrically
    \item \textbf{Pathophysiology:} prolonged abnormal loading strains ligaments, joints, and muscles, causing pain, fatigue, and, over time, degenerative change
\end{itemize}

\section*{Clinical Features}
\begin{itemize}
    \item Neck pain, stiffness, and tension headaches
    \item Upper back and shoulder discomfort
    \item Lower back pain, worse after sitting or standing
    \item Rounded shoulders and a forward head position
    \item Fatigue and a feeling of muscle tightness
    \item Reduced flexibility and range of movement
    \item Pain that eases with movement and posture change
\end{itemize}

\section*{Differential Diagnosis}
\begin{itemize}
    \item Postural pain must be distinguished from other causes of back and neck pain, such as disc problems and arthritis
    \item Headaches may be tension-type or have other causes
    \item Serious causes of back pain, including fracture and infection, are rare but excluded by assessment
    \item Weakness or numbness suggests nerve involvement
    \item A physiotherapist can assess posture and exclude other diagnoses
\end{itemize}

\section*{When to Seek Medical Care}
See a doctor or physiotherapist for persistent pain, pain that wakes you at night, or pain with weakness, numbness, or weight loss.

\textbf{Call 000 (or local emergency services) for:}
\begin{itemize}
    \item Severe back or neck pain after injury
    \item Loss of bladder or bowel control with back pain
    \item Progressive weakness or numbness in the limbs
    \item Back pain with fever
    \item Sudden severe headache or collapse
\end{itemize}

\section*{Diagnosis and Investigations}
\begin{itemize}
    \item \textbf{Postural assessment:} observation of sitting, standing, and movement patterns
    \item \textbf{Clinical examination:} muscle strength, flexibility, and pain provocation
    \item \textbf{Workstation review:} ergonomic assessment of desk, chair, screen, and devices
    \item \textbf{Functional assessment:} how posture affects daily activities
    \item \textbf{Imaging:} rarely needed unless other pathology is suspected
\end{itemize}

\section*{Management}
\begin{itemize}
    \item \textbf{Ergonomic setup:} adjust chair, screen, and keyboard heights; support the lower back
    \item \textbf{Break the habit:} set reminders to change position and take movement breaks every 30--60 minutes
    \item \textbf{Strengthening:} core, gluteal, and upper back exercises correct muscle imbalance
    \item \textbf{Stretching:} chest, hip flexor, and neck stretches relieve tightness
    \item \textbf{Phone and device use:} bring screens to eye level and use voice or larger displays
    \item \textbf{Physiotherapy:} individualised programs and hands-on treatment
    \item \textbf{Activity:} regular walking, swimming, or gym work supports overall posture
    \item \textbf{Awareness:} mindfulness of position during the day
\end{itemize}

\section*{Complications}
\begin{itemize}
    \item Chronic neck, shoulder, and back pain
    \item Tension headaches and jaw clenching
    \item Reduced work performance and productivity
    \item Fatigue and reduced enjoyment of activity
    \item Long-term joint wear and disc degeneration in some people
    \item Poor balance and falls risk in older adults
\end{itemize}

\section*{Prognosis and Outcomes}
Posture improves with consistent effort: most people notice reduced pain and discomfort within weeks of changing habits and starting targeted exercise. The key is sustainability, because posture is maintained by ongoing habits rather than a single fix. Physiotherapy helps those with established pain, and ergonomic improvements prevent recurrence. With attention to positioning, movement, and strength, most people achieve comfortable, healthy posture and freedom from postural pain.

\section*{Prevention and Self-Care}
\begin{itemize}
    \item Set up your workstation ergonomically
    \item Take regular breaks to stand, stretch, and move
    \item Keep screens at eye level and avoid prolonged looking down
    \item Strengthen your core, back, and glutes with regular exercise
    \item Stretch your chest, shoulders, and hips daily
    \item Use supportive chairs and pillows
    \item Carry bags evenly and avoid heavy single-sided loads
    \item Seek physiotherapy advice for persistent postural problems
\end{itemize}

\section*{Sources and Further Reading}
The following reputable sources provide further information for healthcare students and patients seeking reliable, current advice about improving posture.

\begin{itemize}
    \item Healthdirect. \textit{Posture: how to improve it.} https://www.healthdirect.gov.au/posture
    \item Better Health Channel. \textit{Posture and back health.} https://www.betterhealth.vic.gov.au/
    \item Australian Physiotherapy Association. \textit{Posture and pain information.} https://australian.physio/
    \item WorkSafe Australia. \textit{Ergonomic workstation advice.} https://www.safeworkaustralia.gov.au/
    \item NHS. \textit{Posture tips for desk workers.} https://www.nhs.uk/live-well/
    \item Mayo Clinic. \textit{Posture: tips for good posture.} https://www.mayoclinic.org/
\end{itemize}
